\documentclass[11pt,a4paper]{article}
\usepackage[T1]{fontenc}
\usepackage[utf8]{inputenc}
\usepackage[a4paper,margin=2.7cm]{geometry}
\usepackage{amsmath,amssymb,amsthm,mathtools}
\usepackage{microtype}

\newtheorem{theorem}{Věta}
\newtheorem{lemma}[theorem]{Lemma}
\newtheorem{corollary}[theorem]{Důsledek}
\theoremstyle{remark}
\newtheorem{remark}[theorem]{Poznámka}

\newcommand{\FQ}{F_Q}
\newcommand{\HH}{\mathcal H}
\newcommand{\Eop}{\mathsf E}
\newcommand{\Dop}{\mathsf D}
\newcommand{\Sop}{\mathsf S}
\newcommand{\Fop}{\mathsf F}
\newcommand{\Rop}{\mathsf R}

\title{Univerzální renormalizační škálovací pravidlo pro quartic zero-repair}
\author{}
\date{}

\begin{document}
\maketitle

\section{Nastavení a přesný Farkasův lift}

Nechť
\[
P=\left(\frac18,\frac14,\frac14,\frac14,\frac18\right),
\qquad
Q=\left(\frac7{64},\frac{20}{64},\frac{10}{64},\frac{20}{64},\frac7{64}\right).
\]
Pro konečně podporovanou nerostoucí posloupnost
\[
1=a_0\ge a_1\ge a_2\ge\cdots\ge0,
\qquad a_j=0\ \text{pro }j\gg1,
\]
označme symbolem
\[
\mathsf C(a,d):\qquad \sum_{j\ge0}a_j b_j\le d
\]
lineární řez platný pro všechny přípustné rodičovské zákony daného horizontu.
Posloupnost $a$ lze ekvivalentně interpretovat jako survival funkci konečně
podporované náhodné bariéry $Z$:
\[
a_j=\Pr(Z\ge j),\qquad
\Pr(Z=j)=a_j-a_{j+1}.
\]

Položme
\[
\FQ(m):=\Pr_{Y\sim Q}(Y\le m),
\]
s konvencí $\FQ(m)=0$ pro $m<0$ a $\FQ(m)=1$ pro $m\ge4$.

\begin{lemma}[Přesný Farkasův lift]\label{lem:farkas}
Nechť $\mathsf C(a^{(f)},d_f)$ jsou již dokázané řezy na přísně nižších
horizontech. Zvolme reálná $e_x$, nezáporná $u_k$ a nezáporná $v_{xf}$ tak,
aby pro všechna $x\in\{0,1,2,3,4\}$ a $j\ge0$ platilo
\begin{equation}\label{eq:farkas-pointwise}
 e_x+\sum_{1\le k\le x+j}u_k
 +\sum_f v_{xf}a^{(f)}_j\ge0.
\end{equation}
Pak je na následujícím horizontu platný řez
\[
\sum_{j\ge0}A_j b_j\le D,
\]
kde
\begin{align}
A_j&=\sum_{k\ge1}u_k\FQ(k-j-1),\label{eq:Afromu}\\
D&=\sum_xP_xe_x+\sum_ku_k+\sum_{x,f}P_xv_{xf}d_f.\label{eq:Dfromdual}
\end{align}
Je-li $A_0>0$, lze celý řez vydělit $A_0$ a normalizovat na $a_0=1$.
\end{lemma}

\begin{proof}
Pišme $f_{xj}=P_x\Pr(J_x=j)$ pro fan proměnné. Normalizace dětí dává
$\sum_jf_{xj}=P_x$. Stochastické superfan nerovnosti mají tvar
\[
\sum_{x,j:\,x+j\ge k}f_{xj}
\le \Pr(Y+J\ge k),
\qquad Y\sim Q,
\]
a každý starší řez dává
\[
\sum_j a^{(f)}_j f_{xj}\le P_xd_f.
\]
Vynásobíme normalizační rovnosti koeficienty $e_x$, superfan nerovnosti
koeficienty $u_k\ge0$ a starší řezy koeficienty $v_{xf}\ge0$ a vše sečteme.
Koeficient u každého $f_{xj}$ je právě levá strana
\eqref{eq:farkas-pointwise}, a je tedy nezáporný. Proto
\[
0\le
\sum_xP_xe_x+\sum_k u_k\Pr(Y+J\ge k)
+\sum_{x,f}P_xv_{xf}d_f.
\]
Pro pevné $J=j$ platí
\[
\Pr(Y+j\ge k)=1-\FQ(k-j-1),
\]
a tedy
\[
\sum_ku_k\Pr(Y+J\ge k)
=
\sum_ku_k-\sum_jb_j\sum_ku_k\FQ(k-j-1).
\]
Po dosazení dostáváme přesně $\sum_jA_jb_j\le D$.
\end{proof}

\section{Weak--$L^1$ bezpečnostní funkcionál}

Nechť $R_{1/2}$ je kritický weak--$L^1$ kořen se survival funkcí
\[
\Pr(J\ge j)=\frac1{2j},\qquad j\ge1.
\]
Pro bariéru $Z$ se survival funkcí $a$ definujme
\[
\alpha(a):=a_1,
\qquad
h(a):=\mathbb E\frac1{Z+1}
=\sum_{j\ge0}\frac{a_j-a_{j+1}}{j+1}.
\]
Pak
\begin{equation}\label{eq:weak-eval}
\sum_ja_j\Pr(J=j)
=\Pr(J\le Z)
=1-\frac12h(a).
\end{equation}
Bezpečnostní slack řezu $\mathsf C(a,d)$ vůči $C=1/2$ je proto
\begin{equation}\label{eq:slack-general}
s(a,d):=d-1+\frac12h(a).
\end{equation}
Řez je splněn kořenem $R_{1/2}$ právě tehdy, když $s(a,d)\ge0$.

Pro HEAD bariéru
\[
H_{\lambda,d}:\qquad a=(1,\lambda,0,0,\ldots)
\]
platí
\[
h=1-\frac\lambda2,
\qquad
s_H=d-\frac12-\frac\lambda4.
\]
Budeme pracovat s intervalem
\begin{equation}\label{eq:Hcone}
\boxed{
\HH:=\left\{H_{\lambda,d}:\frac18\le\lambda\le\frac7{27},\quad
s_H\ge0\right\}.}
\end{equation}

\section{Čtyři lokální transformace}

\subsection{HEAD reset $\Sop$}

\begin{lemma}[Univerzální one-step škálování]\label{lem:S}
Z libovolného normalizovaného platného řezu $\mathsf C(a,d)$ vznikne jedním
Farkasovým liftem HEAD řez
\begin{equation}\label{eq:Smap}
\boxed{
\Sop(a,d)=H_{\,7/27,\,(8+16d)/27}.}
\end{equation}
Tento vzorec nezávisí na absolutní hloubce.
\end{lemma}

\begin{proof}
V Lematu~\ref{lem:farkas} zvolme
\[
u_2=1,
\qquad
e_0=0,
\qquad e_1=e_2=e_3=e_4=-1,
\qquad v_{1}=1,
\]
a všechny ostatní multiplikátory nulové. Pro dítě $x=1$ je koeficient
v \eqref{eq:farkas-pointwise} roven $-1+a_0=0$ při $j=0$ a $a_j\ge0$
pro $j\ge1$. Pro $x\ge2$ je roven nule a pro $x=0$ je nezáporný.
Z \eqref{eq:Afromu}
\[
A_0=\FQ(1)=\frac{27}{64},
\qquad
A_1=\FQ(0)=\frac7{64},
\qquad
A_j=0\quad(j\ge2).
\]
Po normalizaci tedy $\lambda'=7/27$. Z \eqref{eq:Dfromdual}
\[
D_{\rm raw}=-\frac78+1+\frac14d=\frac18+\frac14d,
\]
a dělení $27/64$ dává $d'=(8+16d)/27$.
\end{proof}

Je-li $H_{\lambda,d}\in\HH$ a
\[
s:=d-\frac12-\frac\lambda4\ge0,
\]
pak pro první dva resety platí
\begin{align}
 s_{\Sop H}
 &=\frac{3+16\lambda+64s}{108}>0,\label{eq:Sslack1}\\
 s_{\Sop^2 H}
 &=\frac{241+256\lambda+1024s}{2916}>0.\label{eq:Sslack2}
\end{align}

\subsection{ESCAPE lift $\Eop$}

Položme
\[
\rho:=\frac\lambda{1-\lambda}
\]
a v Lematu~\ref{lem:farkas} použijme jedinou HEAD závislost
$H_{\lambda,d}$ s multiplikátory
\begin{align*}
u_{m+2}&=\rho^m, &&m=0,1,2,3,4,\\
v_x&=\frac{\rho^{x-1}}{1-\lambda}, &&x=1,2,3,4,\\
e_0&=0,
\qquad
e_x=-\sum_{m=0}^{x-2}\rho^m-v_x, &&x=1,2,3,4,
\end{align*}
se standardní konvencí, že prázdný součet je nula.

\begin{lemma}[Přesný ESCAPE lift]\label{lem:E}
Pro $0<\lambda<1$ výše uvedené multiplikátory splňují
\eqref{eq:farkas-pointwise}. Po normalizaci vznikne bariéra
$G=\Eop(H_{\lambda,d})$ se survival koeficienty
\begin{align}
g_0&=1,\notag\\
g_1&=\frac{24\lambda^4+36\lambda^3-2\lambda^2-\lambda+7}{\Delta},\notag\\
g_2&=\frac{\lambda(2\lambda+1)(20\lambda^2-8\lambda+7)}{\Delta},\notag\\
g_3&=\frac{\lambda^2(17\lambda^2+13\lambda+7)}{\Delta},\label{eq:Eshape}\\
g_4&=\frac{\lambda^3(20\lambda+7)}{\Delta},\notag\\
g_5&=\frac{7\lambda^4}{\Delta},
\qquad g_j=0\quad(j\ge6),\notag
\end{align}
kde
\begin{equation}\label{eq:Delta}
\Delta=47\lambda^4-47\lambda^3+108\lambda^2-71\lambda+27.
\end{equation}
Jeho pravá strana je
\begin{equation}\label{eq:dE}
\boxed{
 d_E(\lambda,d)
 =\frac{8\bigl(\kappa(\lambda)d+\theta(\lambda)\bigr)}{\Delta},}
\end{equation}
kde
\[
\kappa(\lambda)=2-4\lambda+4\lambda^2-\lambda^3,
\qquad
\theta(\lambda)=1-3\lambda+6\lambda^2-2\lambda^3+5\lambda^4.
\]
\end{lemma}

\begin{proof}
Pro $x\ge1$ je výraz v \eqref{eq:farkas-pointwise} při $j=0,1,2$
roven přesně nule. Skutečně používáme identity
\[
\frac1{1-\lambda}=1+\rho,
\qquad
v_x(1-\lambda)=\rho^{x-1}.
\]
Pro $j>2$ může kumulativní součet $u_k$ už jen vzrůst, zatímco HEAD
koeficient je nulový, takže výraz zůstává nezáporný. Pro $x=0$ je tvrzení
okamžité z $u_k\ge0$.

Z \eqref{eq:Afromu} před normalizací
\[
A_j=\sum_{m=0}^4\rho^m\FQ(m+1-j).
\]
Jde o nezápornou nerostoucí posloupnost. Přímé sjednocení jmenovatelů dává
\[
A_0=\frac{\Delta}{64(1-\lambda)^4}>0,
\]
a po vydělení $A_0$ dostáváme \eqref{eq:Eshape}. Dosazení stejných
multiplikátorů do \eqref{eq:Dfromdual} dává po stejné normalizaci
\eqref{eq:dE}.
\end{proof}

Pro další použití zaznamenejme přesný weak--$L^1$ slack ESCAPE stavu. Je-li
$d=\frac12+\frac\lambda4+s$, pak
\begin{equation}\label{eq:sE}
 s_E
 =\frac{8\kappa(\lambda)}{\Delta}s
 +\frac{N_E(\lambda)}{120\Delta},
\end{equation}
kde
\[
N_E(\lambda)=461\lambda^4+174\lambda^3+205\lambda^2-100\lambda+90.
\]
Na intervalu $\lambda\in[1/8,7/27]\subset(0,1/3)$ je
\[
\kappa'(\lambda)=-4+8\lambda-3\lambda^2< -\frac43,
\qquad
\kappa(\lambda)\ge\kappa(1/3)=\frac{29}{27}>0,
\]
a zároveň
\[
N_E(\lambda)\ge 90-100\lambda
\ge 90-\frac{100}{3}=\frac{170}{3}>0.
\]
Tedy $s_E>0$.

\subsection{DELAY lift $\Dop$}

\begin{lemma}[Přesný DELAY lift]\label{lem:D}
Nechť $G=(g,\delta)$ je libovolný normalizovaný bariérový řez a $Z$ jeho
bariéra. Nechť $Y\sim Q$ je nezávislé na $Z$. Potom jedním Farkasovým
liftem vznikne řez $\Dop G$ s bariérou
\begin{equation}\label{eq:delayZ}
\boxed{Z'=Z+4-Y}
\end{equation}
a pravou stranou
\begin{equation}\label{eq:delayd}
\boxed{\delta'=\frac78+\frac18\delta.}
\end{equation}
\end{lemma}

\begin{proof}
Položme $p_z=\Pr(Z=z)=g_z-g_{z+1}$. V Lematu~\ref{lem:farkas} zvolme
\[
u_{5+z}=p_z\quad(z\ge0),
\qquad e_4=-1,
\qquad v_{4,G}=1,
\]
a všechny ostatní multiplikátory nulové. Pro $x=4$ a každý $j\ge0$ platí
\[
-1+\sum_{k\le4+j}u_k+g_j
=-1+(1-g_j)+g_j=0,
\]
zatímco pro $x<4$ jsou koeficienty zjevně nezáporné. Dále
\[
A_m
=\sum_zp_z\FQ(z+4-m)
=\Pr(Z+4-Y\ge m),
\]
což je přesně survival funkce \eqref{eq:delayZ}. Konečně
\[
D=-\frac18+1+\frac18\delta
=\frac78+\frac18\delta.
\]
\end{proof}

Důležitá je nyní specializace na $G=\Dop\Eop(H_{\lambda,d})$.
Označme
\[
\alpha:=\Pr(Z'\ge1),
\qquad
h:=\mathbb E\frac1{Z'+1},
\qquad
\eta:=1-\frac\alpha2-h.
\]
Z \eqref{eq:Eshape} a \eqref{eq:delayZ} dostáváme
\begin{equation}\label{eq:alphaDE}
\alpha
=1-\frac7{64}(1-g_1)
=\frac{N_\alpha(\lambda)}{64\Delta},
\end{equation}
kde
\[
N_\alpha(\lambda)
=2847\lambda^4-2427\lambda^3+6142\lambda^2-4054\lambda+1588,
\]
a
\begin{equation}\label{eq:hDE}
h
=\frac{1992081\lambda^4-4128338\lambda^3+7213185\lambda^2
-4834200\lambda+1661940}{161280\Delta}.
\end{equation}
Dvě rozhodující veličiny se zjednoduší na
\begin{align}
16\eta-5\alpha+2
&=\frac{N_\Phi(\lambda)}{20160\Delta},\label{eq:phase-margin}\\
s_{DE}
&=\frac{\kappa(\lambda)}{\Delta}s
+\frac{N_D(\lambda)}{322560\Delta},\label{eq:DEslack}
\end{align}
kde
\begin{align*}
N_\Phi(\lambda)
&=1412733\lambda^4+1139881\lambda^3-386820\lambda^2
+505050\lambda-28980,\\
N_D(\lambda)
&=1629201\lambda^4-2717138\lambda^3+5116545\lambda^2
-3423000\lambda+1218420.
\end{align*}

\begin{lemma}[Fusion-ready vlastnosti delayed ESCAPE stavu]\label{lem:DEsafe}
Pro každý $H_{\lambda,d}\in\HH$ platí pro
$G=\Dop\Eop(H_{\lambda,d})$
\begin{equation}\label{eq:DEproperties}
\boxed{
s_{DE}>0,
\qquad
16\eta-5\alpha+2>0,
\qquad
\frac{57}{64}\le\alpha\le1.}
\end{equation}
\end{lemma}

\begin{proof}
Již víme, že $\Delta>0$ a $\kappa>0$. Pro $N_D$ použijeme čistě
algebraický certifikát pozitivity. Položme $t=3\lambda$. Na $0\le\lambda\le1/3$
platí $t\in[0,1]$ a
\[
N_D(\lambda)
=\sum_{k=0}^4 c_k\binom4k t^k(1-t)^{4-k},
\]
kde
\[
(c_0,c_1,c_2,c_3,c_4)
=\left(
1218420,\ 933170,\ \frac{4456025}{6},\
\frac{16787623}{27},\ \frac{15265904}{27}
\right).
\]
Všechny Bernsteinovy koeficienty jsou kladné, tedy $N_D>0$ a
\eqref{eq:DEslack} dává $s_{DE}>0$.

Pro $N_\Phi$ máme na $[1/8,7/27]$
\[
N_\Phi'(\lambda)
=5650932\lambda^3+3419643\lambda^2-773640\lambda+505050
\ge \frac{913430}{3}>0.
\]
Proto
\[
N_\Phi(\lambda)
\ge N_\Phi(1/8)
=\frac{125658821}{4096}>0,
\]
a \eqref{eq:phase-margin} je kladná.

Konečně z \eqref{eq:delayZ} nastane $Z'=0$ pouze tehdy, když současně
$Z=0$ a $Y=4$. Proto
\[
\alpha=1-\frac7{64}\Pr(Z=0)
\ge1-\frac7{64}=\frac{57}{64},
\]
a samozřejmě $\alpha\le1$.
\end{proof}

\subsection{FUSION lift $\Fop$}

\begin{lemma}[Přesná fusion formule]\label{lem:F}
Nechť $G=(g,\delta)$ je bariérový řez s
\[
\alpha=g_1,
\qquad
h=\mathbb E\frac1{Z+1},
\qquad
\eta=1-\frac\alpha2-h,
\qquad
s_G=\delta-1+\frac h2,
\]
a nechť $H_{\beta,d}$ je HEAD řez se slackem
\[
s_H=d-\frac12-\frac\beta4.
\]
Pak existuje přesný one-step Farkasův lift
\[
\Fop(G,H_{\beta,d})=H_{\lambda',d'},
\]
kde
\begin{equation}\label{eq:fusion-map}
\boxed{
\lambda'=\frac{7\alpha}{7+20\alpha},
\qquad
d'=\frac{8\delta+16\alpha d}{7+20\alpha}.}
\end{equation}
Jeho weak--$L^1$ slack je přesně
\begin{equation}\label{eq:fusion-slack}
\boxed{
s'
=\frac{32s_G+64\alpha s_H+16\eta+2-7\alpha+16\alpha\beta}
{4(7+20\alpha)}.}
\end{equation}
Zvláště, jestliže
\begin{equation}\label{eq:fusion-sufficient}
16\eta-5\alpha+2\ge0,
\qquad
\beta\ge\frac18,
\qquad
s_G,s_H\ge0,
\end{equation}
pak $s'\ge0$.
\end{lemma}

\begin{proof}
Zvolme libovolné $T>0$ a v Lematu~\ref{lem:farkas} položme
\[
u_1=T(1-\alpha),
\qquad
u_2=T\alpha,
\qquad
e_x=-T\quad(0\le x\le4),
\]
\[
v_{0,G}=T,
\qquad
v_{1,H}=T\alpha,
\]
a ostatní $v$ rovny nule. Pro $x=0$ se rovnost v
\eqref{eq:farkas-pointwise} saturuje při $j=0,1$ a pro $j\ge2$ zůstává
$Tg_j\ge0$. Pro $x=1$ se saturuje při $j=0$ a při $j=1$ zbývá
$T\alpha\beta\ge0$. Pro $x\ge2$ se $u_1+u_2=T$ přesně vyruší s $e_x=-T$.

Z \eqref{eq:Afromu}
\[
A_0=\frac{T(7+20\alpha)}{64},
\qquad
A_1=\frac{7T\alpha}{64},
\qquad
A_j=0\quad(j\ge2),
\]
což po normalizaci dává první formuli v \eqref{eq:fusion-map}. V
\eqref{eq:Dfromdual} se členy $\sum_xP_xe_x=-T$ a $\sum_ku_k=T$ zruší;
zůstane
\[
D_{\rm raw}=T\left(\frac\delta8+\frac{\alpha d}{4}\right),
\]
což dává druhou formuli v \eqref{eq:fusion-map}.

Dosadíme
\[
\delta=1-\frac h2+s_G,
\qquad
d=\frac12+\frac\beta4+s_H,
\qquad
h=1-\frac\alpha2-\eta
\]
do $d'-1/2-\lambda'/4$. Po přesném zkrácení vznikne
\eqref{eq:fusion-slack}. Poslední část plyne z identity
\[
16\eta+2-7\alpha+16\alpha\beta
=(16\eta-5\alpha+2)+2\alpha(8\beta-1)\ge0.
\]
\end{proof}

\section{Hlavní škálovací věta}

Definujme tříkrokový renormalizační operátor na HEAD řezech
\begin{equation}\label{eq:Rdef}
\boxed{
\Rop(H):=
\Fop\bigl(\Dop\Eop(H),\,\Sop^2(H)\bigr).}
\end{equation}
První argument fusion i druhý argument fusion mají o dva horizonty méně než
výsledný řez, takže konstrukce je prefixově kauzální a neobsahuje kruhové
odůvodnění.

\begin{theorem}[Univerzální renormalizační škálovací pravidlo]\label{thm:scaling}
Pro každý horizont, na němž je HEAD řez $H_{\lambda,d}\in\HH$ již dokázán,
je o tři úrovně výše dokázán HEAD řez
\[
\Rop(H_{\lambda,d})=H_{\Lambda(\lambda),D(\lambda,d)}
\]
s
\begin{equation}\label{eq:Lambda}
\boxed{
\Lambda(\lambda)
=\frac{7\alpha(\lambda)}{7+20\alpha(\lambda)},}
\end{equation}
kde
\begin{equation}\label{eq:alpha-explicit}
\alpha(\lambda)
=\frac{2847\lambda^4-2427\lambda^3+6142\lambda^2-4054\lambda+1588}
{64\Delta(\lambda)},
\end{equation}
a
\begin{equation}\label{eq:Dmap}
\boxed{
D(\lambda,d)
=
\frac{8\delta(\lambda,d)+16\alpha(\lambda)d_2(d)}
{7+20\alpha(\lambda)},}
\end{equation}
kde
\[
\delta(\lambda,d)=\frac78+\frac18d_E(\lambda,d),
\qquad
d_2(d)=\frac{344+256d}{729}.
\]
Navíc
\begin{equation}\label{eq:inclusion}
\boxed{\Rop(\HH)\subseteq\HH.}
\end{equation}
Tvrzení je nezávislé na absolutní hloubce; hloubka vstupuje pouze přes
požadavek, že použité řezy již byly dokázány na přísně nižších horizontech.
\end{theorem}

\begin{proof}
Vezměme $H_{\lambda,d}\in\HH$ a označme jeho slack $s\ge0$.
Lemma~\ref{lem:DEsafe} říká, že
\[
G:=\Dop\Eop(H_{\lambda,d})
\]
má $s_G>0$, splňuje
\[
16\eta_G-5\alpha_G+2>0
\]
a zároveň $57/64\le\alpha_G\le1$.

Na druhé větvi použijeme dvakrát Lemma~\ref{lem:S}. Z
\eqref{eq:Sslack2} dostáváme
\[
\Sop^2(H_{\lambda,d})=H_{7/27,d_2(d)},
\qquad s_{\Sop^2H}>0.
\]
V Lematu~\ref{lem:F} tedy můžeme vzít $\beta=7/27>1/8$; všechny podmínky
\eqref{eq:fusion-sufficient} jsou splněny. Výstupní HEAD řez má proto
$s'\ge0$.

Zbývá ověřit interval pro nový shape parametr. Funkce
\[
f(\alpha)=\frac{7\alpha}{7+20\alpha}
\]
je rostoucí, protože
\[
f'(\alpha)=\frac{49}{(7+20\alpha)^2}>0.
\]
Proto z $57/64\le\alpha\le1$ plyne
\[
\frac{399}{1588}
=f\!\left(\frac{57}{64}\right)
\le\Lambda(\lambda)
\le f(1)=\frac7{27}.
\]
Navíc $399/1588>1/4>1/8$. Tedy
\[
\frac18<\Lambda(\lambda)\le\frac7{27},
\]
a spolu s $s'\ge0$ dostáváme
$\Rop(H_{\lambda,d})\in\HH$. To dokazuje \eqref{eq:inclusion}.
\end{proof}

\begin{corollary}[Přesná reprodukce zlomu v šestém horizontu]\label{cor:n6}
Pro HEAD řez
\[
\lambda=\frac7{27},
\qquad
d=\frac{5144}{6561}
\]
dává jednotlivý renormalizační cyklus přesně
\begin{align*}
d_E&=\frac{202152433}{246095334},\\
\delta&=\frac{1924819771}{1968762672},\\
d_2&=\frac{3573848}{4782969},\\
\Lambda(\lambda)&=\frac{855963003}{3353021828},\\
D(\lambda,d)&=\frac{986801982722048}{1336449954970611}.
\end{align*}
Odpovídající weak--$L^1$ dolní mez je
\[
\frac{1-D(\lambda,d)}{1-\Lambda(\lambda)/2}
=\frac{2797183777988504}{9326918136932919}
=0.299904398958\ldots,
\]
tedy přesně hodnota, která nahrazuje jednoduchou HEAD rekurenci v šestém
kroku.
\end{corollary}

\begin{remark}[Přesný rozsah věty]
Věta~\ref{thm:scaling} nedokazuje, že libovolný Farkasův separator leží v
HEAD/PHASE gramatice, ani sama o sobě nedokazuje $C_n\le1/2$ pro všechna
$n$. Dokazuje přesně jednu load-bearing skutečnost: konkrétní
renormalizační grammar
\[
H\longmapsto \Fop\bigl(\Dop\Eop(H),\Sop^2(H)\bigr)
\]
je uzavřená na $\HH$ a její pravidla jsou nezávislá na hloubce. K úplnému
all-horizon důkazu zbývá completeness/dominance lemma, které ukáže, že žádný
jiný rekurzivně vznikající separator nemůže weak--$L^1$ kořen $C=1/2$
vyloučit.
\end{remark}

\begin{remark}[Důležitá kontrolní korekce]
Není pravda, že DELAY operátor $\Dop$ zachovává podmínku
$16\eta-5\alpha+2\ge0$ na celém abstraktním PHASE kuželu. Výše se nic
takového nepoužívá. Dokázána je pouze potřebná specializace
\[
\Dop\Eop(\HH)\subseteq
\{G:s_G\ge0,\ 16\eta_G-5\alpha_G+2\ge0\},
\]
což je přesně obsah Lematu~\ref{lem:DEsafe}. Tato restrikce je nutná:
pro deterministickou bariéru $Z\equiv0$ je vstupní hodnota
$16\eta-5\alpha+2=2$, ale po DELAY kroku vychází
$16\eta'-5\alpha'+2=-251/960<0$. Tím se odstraňuje dřívější příliš silná
formulace a důkaz je uzavřen bez tohoto neplatného kroku.
\end{remark}

\end{document}
